\chapter{A New Fluorescence Energy Scale}
\label{chap:xy-scanner}

The analysis in the subsequent chapters heavily relies on information measured 
by the \WCD, \SSD, and \UMD. With the inauguration of the latter two detectors 
ocurring roughly ten years after the last official measurement of \FD pixel 
calibration constants \cite{dorofeevFDAbsoluteCalibration2013}, the 
cross-calibration of the reconstructed \UMD/\SSD energies (Ref.\ 
\cref{sssec:sd-sd-reconstruction-energy-estimator}) depend on measurement data 
that are more than a decade old. 

While the pixel calibration constants are in principle updated using a nightly 
relative calibration (Ref.~\cref{sssec:relative-calibration}) it is crucial to 
verify the integrity of the relative tracking and the employed energy 
calibration.

The XY-Scanner is a novel end-to-end calibration system that is used to estimate
the \FD pixel calibration constants with an accuracy of \todo{uncertainty}. It 
provides an excellent method of verifying the validity of the currently employed
pixel calibration constants, and the \FD energy scale by extension. The 
measurement system as well as underlying analysis is presented in detail in 
\cite{schaferXYScannerVersatileMethod2023}, and summarized in 
\cref{sec:xy-scanner}.

\cref{sec:quality-assurance} details extensive quality checks that have been
developed in this work, and that are used to assure the reproducibility of XY 
measurements and establish the trustworthiness of the gathered data.

After ensuring the quality of the novel calibration method, the results of 
various XY measurement campaigns are compared to current pixel calibration 
constants can. \cref{sec:telescope-aging} details this comparison and details
possible reasons for discrepancies. An updated relative calibration tracking 
algorithm is equally given.

Finally, an update to the \FD energy scale is presented in 
\cref{sec:fd-energy-scale}.

\section{The XY-Scanner}
\label{sec:xy-scanner}

\begin{figure}[t]
  \centering
  \subfloat[]{\includegraphics[width=0.47\textwidth]{xy-scanner/xy_light_source.jpg}
  \label{fig:xy-light-source}
  }\hspace{0.2cm}
  \subfloat[]{\includegraphics[width=0.47\textwidth]{xy-scanner/xy_pattern.png}
  \label{fig:xy-flash-pattern}
  }
  \caption[]{\subref{fig:xy-light-source} The XY calibration light source 
  mounted in a test setup in Karlsruhe. Photo from 
  \cite{schaferXYScannerVersatileMethod2023} \subref{fig:xy-flash-pattern} The
  XY scanner, consisting of light source, horizontal and vertical axles, as well
  as the electronics box on the top right, mounted in front of an \FD telescope.
  Photo from \cite{schaferXYScannerVersatileMethod2023}. Flasher positions are 
  superimposed as grey hexagons.}
  \label{fig:xy-hardware}
\end{figure}

\fulllift

\subsection{Light Source Calibration}
\label{ssec:xy-absolute-calibration}

The spirit of the \FD calibration procedure relies on the \FD-camera being
exposed to a light source of known radiance. This allows to tie the signal
recorded in the pixel-\PMTs to an incident count of photons, which, during \DAQ 
holds information about the energy of an \EAS (Ref.\ 
\cref{sssec:fd-reconstruction-energy-estimator}).

Different facilities have been set up to measure the radiance of the XY \LED, 
shown in \cref{fig:xy-light-source} and hereafter referred to as
\CLS, independently \cite{heibgesPhotometricCharacterizationLight2020, 
schaferXYScannerVersatileMethod2023, vaculaVyvojKalibracnihoSvetelneho2024}. 
While precise implementation details differ, all test banks rely on measuring a 
distance-dependent flux, $\Phi_\mathrm{PD}(r)$, of photons with wavelength 
$\lambda$ on a \PD illuminated by the \CLS. This flux follows, to first order,
an inverse-square law, supplemented by higher orders that consider the view 
factor $F_{\mathrm{PD},\,\mathrm{CLS}}$ that depends on the aperture radii
$R_\mathrm{PD}$, $R_\mathrm{CLS}$ of the \PD- and \CLS aperture respectively.

\begin{equation}
\label{eq:source-radiance}
F_{\mathrm{PD},\,\mathrm{CLS}} = \frac{1}{2}\,\left[ 1 + \left(\frac{r-r_0}{R_\mathrm{CLS}}\right)^2 + \left(\frac{R_\mathrm{PD}}{R_\mathrm{CLS}}\right)^2 - \sqrt{1 + \left(\frac{r-r_0}{R_\mathrm{CLS}}\right)^2 - 3\left(\frac{R_\mathrm{PD}}{R_\mathrm{CLS}}\right)^2} \right] \\
\end{equation}

The diode at distance $r-r_0$ from the \CLS consequently measures a flux 

\begin{equation}
\Phi_\mathrm{PD}(r) = \frac{hc}{\lambda}\,\tilde{S}_\mathrm{setup}\,L_\mathrm{CLS},
\end{equation}

where $\tilde{S}_\mathrm{setup}$ in \cref{eq:source-radiance} gives the response
function of the \PD.\footnote{Which will in general depend on the entire test
setup, not just the \PD.} By controlling the laboratory environment precisely, 
the source radiance can be determined with an accuracy of 
$\lessapprox\SI{3}{\percent}$. \cref{fig:abs-calib-measurements} and 
\todo{abscalpaper} detail the agreement of the cross-calibration between 
different measurement facilities.

\begin{figure}[t]
  \centering
  \subfloat[]{\includegraphics[height=0.36\textwidth]{xy-scanner/abs_calib_jlo.png}
  \label{fig:abs-calib-jlo}
  }\hspace{0.2cm}
  \subfloat[]{\includegraphics[height=0.36\textwidth]{xy-scanner/abs_calib_kit.png}
  \label{fig:abs-calib-kit}
  }
  \caption[]{The signal response versus distance of a \PD is used to estimate 
  the source radiance of the \CLS. Shown are results for the same \CLS tested 
  \subref{fig:abs-calib-jlo} in Olomouc, Czech Republic, and 
  \subref{fig:abs-calib-kit} in Karlsruhe, Germany.}
  \label{fig:abs-calib-measurements}
\end{figure}

\subsection{Telescope Calibration}
\label{ssec:fd-xy-calibration}

The precisely determined radiance of the \CLS is then used to connect the \FD 
response after a \CLS light flash to a corresponding calibration constant.

For this purpose the XY lightsource is installed in front of the \FD aperture on
two motorized linear stages, as shown in \cref{fig:xy-flash-pattern}. Over the 
course of a measurement the light source tracks a path consisting of 
$n_\mathrm{pos}>1500$ unique positions distributed uniformly on a hexagonal grid
that covers the entire telescope frame. At each position the \CLS emits a flash 
of light directed at the \FD-camera. The detector response is recorded for each 
position.

The measurement process takes on average \SI{40}{\minute} without interruptions.
To check the stability of the \FD \DAQ readout chain, two Calibration A 
measurements with an open \FD shutter are performed immediately before and after
the XY run.

A complete XY measurement thus consist of $n_\mathrm{pos}$ \FD events, each 
containing the \ADC-traces digitized at the time of the \CLS flash for all 440 
pixels of the \FD-camera. The calibration constant for a given pixel-\PMT is 
then calculated from this information according to the steps listed below.

\vspace{1cm}
\begin{enumerate}
	\item The baseline $b_i$ is calculated per \PMT and subtracted from all 
	traces $t_i$.
	\item A pulse finding algorithm 
	(Ref.~\cite[][Section 7.1.1]{schaferXYScannerVersatileMethod2023}) 
	determines the trace interval $[k_\mathrm{start}, k_\mathrm{stop}]$ that
	contains the light pulse from the calibration light source.
	\item The integral of the pulse is calculated, this value is summed for 
	all events in each pixel to yield a total signal per pixel in units of 
	\ADC.
	\item The total number of photons received by each pixel-\PMT is 
	estimated via the known source radiance. The non-uniform camera 
	geometry and reflectivity are taken into account via simulation based 
	correction factors $F_{i,\,\mathrm{CLS}}$.
	\item The total signal in every \PMT is divided by the total count of 
	incident photons to calculate the pixel calibration constant in units of
	$\ADC/\gamma$.
\end{enumerate}

In summary, the calibration constant $C^{\mathrm{pixel},\,i}_\mathrm{abs.}$ for 
pixel $i$ is calculated as 

\begin{equation}
	\label{eq:xy-scanner-calib-calculation}
	C^{\mathrm{pixel},\,i}_\mathrm{abs.} = \sum\limits_{j=1}^{n_\mathrm{pos}} \left[ \sum\limits_{k = k_\mathrm{start}}^{k_\mathrm{stop}} t_{i,\,k} - b_i \right]\,\bigg/\,n_\mathrm{pos}L_\mathrm{CLS}\,F_{i,\,\mathrm{CLS}}.
\end{equation}

The different systematic and statistical contributions intrinsic to this process
are listed in \cref{tab:xy-systematics}. The total uncertainty, given as 
$\sqrt{\sum\sigma_i^2}$, amounts to \todo{uncertainty}.

\begingroup
\renewcommand{\arraystretch}{1.3}
\begin{table}
	\begin{center}
	\caption{Uncertainties of pixel calibration constants measured by the 
	XY-Scanner. Comparisons to the Drum calibration, where possible, are 
	taken from \cite{brackAugerFluorescenceDetector2013, 
	crissUpdateRecentAbsolute2010, dorofeevFDAbsoluteCalibration}.}
	\begin{tabular*}{\textwidth}{w{l}{0.35\textwidth}w{r}{0.35\textwidth}|w{c}{0.1\textwidth}w{c}{0.1\textwidth}}
  	\toprule
	 & & XY & Drum \\
	\hline
	Contribution & details & \multicolumn{2}{c}{$\sigma_i$} \\
	\hline
	Radiance of light source & & & \SI{2.1}{\percent} \\
	Aperture coverage factor & 
	\cite[][p.~127]{schaferXYScannerVersatileMethod2023} &
	\SI{1.3}{\percent} & \SI{2}{\percent} \\
	Directional emission & 
	\cite[][p.~83]{schaferXYScannerVersatileMethod2023} & 
	\SI{1}{\percent} & $-$ \\
	Back reflection on camera & 
	\cite[][p.~47]{schaferXYScannerVersatileMethod2023} & 
	\SI{0.5}{\percent} & \SI{2}{\percent} \\
	Signal in the pixel-\PMTs &
	\cite[][p.~114]{schaferXYScannerVersatileMethod2023} & 
	\SI{0.2}{\percent} & \SI{2}{\percent} \\
	\hline
	Statistical uncertainty & 
	\cite[][p.~116]{schaferXYScannerVersatileMethod2023} & 
	\SI{2.1}{\percent} & $-$ \\
	\hline
	Unaccounted uncertainty budget & & & $\approx\SI{9}{\percent}$ \\
	\hline
	\textbf{Total} & & & \textbf{\SI{9.9}{\percent}} \\ 
  \bottomrule
	\end{tabular*}
	\label{tab:xy-systematics}
	\end{center}
\end{table}
\endgroup


\section{Quality Assurance}
\label{sec:quality-assurance}

By calculating \FD pixel calibration constants as a sum over many individual 
light flashes according to \cref{eq:xy-scanner-calib-calculation} in 
\cref{ssec:fd-xy-calibration}, one implicitly requires several assumptions to be
true, namely,

\vspace{0.4cm}
\begingroup
\renewcommand\arraystretch{1.3}
\begin{itemize}
\item the source radiance $L_\mathrm{CLS}$ remains constant during the 
measurement.
\item the \CLS flashes cover the complete filter, and the aperture is 
unoccluded.
\item the sensitivity of the \FD pixel-\PMTs is constant during the measurement.
\end{itemize}
\endgroup
\vspace{0.4cm}

Thus, in order to ensure that the XY scanner delivers accurate and precise pixel
calibration constants one must verify that these assumptions hold. This can be
done by establishing a checklist of tests that must be executed and that have to
pass before accepting the results of XY measurement runs.

The following sections detail what these tests are, and how failure states are 
defined based on outlier identification and cross-correlation of available XY
data.

\subsection{Preselection of XY Runs}

To date, six measurement campaigns have been conducted that have probed the
calibration constants in all 27 telescopes at least once. With the XY scanner 
actively being developed during the early phases of measuring, many different 
run configurations have been tested throughout the different campaigns. 

A long-term strategy has emerged from the prototyping phase that compromises 
between ease-of-use, long-term feasibility, and systematic uncertainty 
\cite{filipQualityAssuranceXY2024}. The run configuration of the XY scanner 
measurements are defined as:\footnote{Details on the nomenclature are found in 
\cite{schaferXYScannerVersatileMethod2023}.}

\vspace{0.4cm}
\begin{itemize}
\item Use a type \textbf{3P-GPS-053-S} integrating sphere sporting a modified 
baffle to improve the Lambertian diffuse emission profile of the XY light 
source.

\item Install \textbf{Generation 2 LED head} in the \IS, which supports three
\LEDs of type UVLED365-110E \cite{alldatasheet.comUVLED365110EDownload}. A thin 
diffusive material that covers the \LEDs, as well as a shielding baffle improve
the emission profile.

\item Employ a \textbf{stepsize of 6 cm}. The step size between individual 
flasher positions directly influences the coverage factor of the telescope 
aperture. The smaller the step size, the lower the systematic uncertainty of the
calibration constants

\item Flash the \LED using a \textbf{current of 15.9 mA}, which compromises 
between \FD signal strength (and thus accuracy) and long-term operation of the 
\LED.

\end{itemize}
\vspace{0.4cm}

All data shown in the upcoming quality assurance analysis has been preselected
to adhere to these standards. A list of used measurements can be found in
\cref{tab:xy-runlist}.

\halflift
\subsection{Light Source Stability}

The XY \LED is located inside the \CLS alongside a \PD that logs the intensity 
of each flash. While the \PD is not precise enough to track the source radiance
exactly, it is sufficient enough for monitoring the nominal \LED operation. 
Moreover, temperature sensors are installed within the sphere to monitor the 
climate controlled environment of the \CLS. 

\subsubsection{Temperature}
\label{sssec:xy-qa-temperature}

The emission characteristics of the \CLS \LED are temperature dependent 
\cite{alldatasheet.comUVLED365110EDownload}. For this reason, the inside of the
\IS is kept at a constant operating temperature of $\SI{30}{\celsius}$. External
factors such as wind gusts can cause temporary deviations from the nominal 
temperature, in an extreme case, this influences the total emitted signal, and
pixel calibration constants as a consequence.

In the following, we estimate the influence of temperature on the emissivity. We
do this by examining both the signal seen by the \PD, as well as the much more 
accurate \ADC integral recorded by the \FD camera, normalized to the source
radiance at the time of measurement.

The \PD signal measures the \LED intensity directly and can be taken as is. For 
the \FD signal, to isolate effects of temperature, only flashes at the center of
the aperture are compared to one another. Moreover, to account for varying \FD 
camera efficiency the integral signal is normalized by the mean camera signal 
during the open shutter Cal A measurements before and after the corresponding 
run.

\begin{figure}[t]
  \centering
  \includegraphics[width=1.0\textwidth]{xy-scanner/led_temperature_stability.png}
  \caption{The influence of the \LED temperature on the flash intensity as seen
  by the \FD camera (blue) and the internal \PD (red). The dashed lines give 
  linear fits to the mean signal per temperature for temperatures with $>10$ 
  recorded flashes.}
  \label{fig:temperature-stability}
\end{figure}

\cref{fig:temperature-stability} shows the distribution of recorded signals for
the \FD and the \PD. No clear trend is identifiable in both the \PD response, as
well as the \FD camera signal.

Expecting in theory a positive correlation between temperature and \LED 
emissivity in the considered temperature ranges 
(Ref.~\cite[p.~2, lower left]{alldatasheet.comUVLED365110EDownload}), we 
construct quality cuts by forming the very cautious assumption that flashes 
outside the temperature range from $\SIrange{29}{31}{\celsius}$ deviate from the
mean signal by $\approx\SI{5}{\percent}$. 

With $n_\mathrm{pos} \gg 1000$ flashes per run, a single flash outside of this
2-degree range impacts the total signal by less $1\permil$. Limiting the 
allowed magnitude of the temperature dependency on the emitted signal to 
$<\SI{0.2}{\percent}$ (Ref.~\cref{tab:xy-systematics}), we accept an XY run if 
it has $\leq4$ flashes where 
$T_\mathrm{LED}\notin[\SI{29}{\celsius},\,\SI{31}{\celsius}]$.

While the aforementioned assumption of a $\SI{5}{\percent}$ bias is very 
conservative, the effect on data is marginal. Most runs have all data points 
within $\SIrange{29}{31}{\celsius}$. There exists one run with a single flash at
$\SI{28.5}{\celsius}$, and two (rejected) runs with 11 and 19 bad temperatures 
respectively, clearly marked as outliers in this analysis.

\subsubsection{Temporal Drifts}
\label{sssec:xy-qa-led-drift}

The pulse-to-pulse stability of the \CLS during radiance measurements is 
\SI{2}{\percent} \cite[p.~100]{schaferXYScannerVersatileMethod2023}. We can 
verify that this assumption holds during XY scans as well by examining the data
from the \PD mounted inside the integrating sphere. 

Examining the same data as in \cref{sssec:xy-qa-temperature}, we identify 
possible drifts in the source radiance by fitting a first order polynomial to 
the \PD signal for each run. In doing this, we need to account for the inherent 
variance in the \PD signal, which is far greater than the systematic uncertainty
of the remaining components. Examining all runs thus far, $\sigma_\mathrm{PD}$ 
seems to be well described by a normal distribution 
$\mathcal{N}(\mu = \SI{0.583}{\ADC\per\giga\photon},\, 
\sigma^2=\SI{0.009}{\square\ADC\per\giga\square\photon})$ 
(Ref.~\cref{fig:sigma-pd}).

With the resulting linear model $\hat{S}_\mathrm{PD}$ that approximates the 
measured signal $S_\mathrm{PD}$, it becomes trivial to check the minimum 
pulse-to-pulse stability by calculating the percentual difference in $\hat{S}$ 
between the first and last emitted flash of the XY run.

\cref{fig:pd-drift} shows the distribution of this percentual difference. The
\PD, and by extension the \CLS \LED operate stable during \DAQ. All runs pass 
this quality test. 

In practice, we are unable to distinguish a drift in the \LED from a drift in 
the \PD, meaning there exists the possibility to discard "good" data due to a 
"bad" \PD. This makes the above quality cut a cautious (but for the time being) 
correct choice.

\begin{figure}[t]
  \centering
  \subfloat[]{\includegraphics[height=0.38\textwidth]{xy-scanner/sigma_pd.png}
  \label{fig:sigma-pd}
  }\hspace{0.2cm}
  \subfloat[]{\includegraphics[height=0.38\textwidth]{xy-scanner/pd_drift.png}
  \label{fig:pd-drift}
  }
  \caption[]{\subref{fig:sigma-pd} Noise level of the \CLS \PD, normalized to
  the respective $L_\mathrm{CLS}$. \subref{fig:pd-drift} Difference in signal 
  strength between the first and last flash extrapolated from a linear fit. The 
  red histogram marks runs that have already been rejected in \cref{sssec:xy-qa-temperature}.}
  \label{fig:pd-noise-level-and-drift}
\end{figure}

\subsubsection{Interruptions}

Next, we estimate the influence of \DAQ interruptions on the \CLS stability.
These interruptions may occur if the XY scanner encounters hardware/software
problems during a measurement, and were especially frequent prior to November 
2023 \cite{filipFixfilespy2023}.

For the sake of this discussion, an interruption is classified as the scenario
in which more than $\SI{20}{\second}$ pass between consecutive flashes (Ref.)

\subsection{Measurement Setup}



\subsection{Camera Stability}
\subsection{Reproducibility}

\section{Telescope Ageing}
\label{sec:telescope-aging}

% address \cref{tab:xy-systematics} unaccounted uncertainty budget and why

\subsection{Comparison of Drum and XY}
\subsection{The Night Sky Background}
\subsection{Effects Of Mirror Cleaning}

\section{FD Energy Scale}
\label{sec:fd-energy-scale}


\section{Summary}
